% ==============================================================================
% 03-functional.tex — Раздел 2.2. Функциональные требования (сокращённо)
% ==============================================================================

\section{Функциональные требования}
\label{sec:functional}

Перечень \abbr{функциональных требований}{ФТ} в~соответствии с~практиками инженерии требований~\cite{src:8} приведён в~таблице~\ref{tab:fr}.

\begin{table}[H]
  \caption{Функциональные требования}
  \label{tab:fr}
  \small
  \begin{tabularx}{\textwidth}{|c|p{3.2cm}|X|}
    \hline
    \textbf{ID} & \textbf{Название} & \textbf{Описание} \\
    \hline
    FR-01 & Карточка бизнеса & Создание и ведение карточки (название, категория, адрес, телефон, описание); роли владелец или администратор \\
    \hline
    FR-02 & Расписание & Часы работы по дням и особым датам; эталон для синхронизации \\
    \hline
    FR-03 & Подключение платформ & Привязка VK, Telegram, Яндекс.Бизнес; хранение токенов в~зашифрованном виде \\
    \hline
    FR-04 & Отключение платформ & Удаление или отключение интеграции; роли владелец или администратор \\
    \hline
    FR-05 & Консистентность данных & Эталонные и фактические значения по каналам; обнаружение расхождений, задачи синхронизации \\
    \hline
    FR-06 & Оркестрация агентов & Запрос на естественном языке → языковая модель → вызовы инструментов → агенты → результаты и потоковая передача ответа \\
    \hline
    FR-07 & Входящие события & Регистрация сообщений, отзывов, комментариев; статусы обработки \\
    \hline
    FR-08 & Журналирование & История изменений, статусы задач, трассируемость действий \\
    \hline
    FR-09 & Регистрация & Электронная почта, пароль (не менее 8 символов); создание пользователя и бизнеса по умолчанию \\
    \hline
    FR-10 & Аутентификация & Вход по почте и паролю; выдача токенов доступа и обновления \\
    \hline
    FR-11 & Разграничение прав & Роли по таблице~\ref{tab:roles}; проверка при каждом запросе \\
    \hline
  \end{tabularx}
\end{table}
